\documentclass[../../../include/open-logic-section]{subfiles}

\begin{document}

\olfileid{sth}{cardinals}{hp}
\olsection{附录：休谟原则}

在 \olref[cp]{sec} 中，我们描述了康托原则。它是：
\begin{align*}
	\card{A} = \card{B} & \text{ 当且仅当 } A \approx B.
\intertext{这与现在所谓的\emph{休谟原则}非常相似，后者说的是：}
	\fregenum{x} {F(x)} = \fregenum{x}{G(x)} & \text{ 当且仅当 } F \sim G
\end{align*}
其中‘$F \sim G$’是“$F$ 和 $G$ 在数量上恰好一样多”的缩写，即 $F$ 与 $G$ 之间可以建立双射，亦即：
\begin{align*}
	\exists R(&\forall v\forall y(Rvy \lif (Fv \land Gy)) \land {}\\
		&\forall v(Fv \lif \lexists![y][Rvy]) \land {}\\
		&\forall y(Gy \lif \lexists![v][Rvy]))
\end{align*}
但休谟原则与康托原则之间存在类型差异。在康托原则的陈述中，变元“$A$”和“$B$”是代表\emph{集合}的一阶项。在休谟原则的陈述中，“$F$”、“$G$”和“$R$”\emph{不是}一阶项；相反，它们处于\emph{谓词位置}。（也许它们代表\emph{性质}。）因此，我们可以用自然语言将休谟原则表述为：$F$ 的数量等于 $G$ 的数量，当且仅当 $F$ 与 $G$ 之间存在双射。这之所以被称为\emph{休谟原则}，是因为休谟曾写道：

\begin{quote}
  当两个数以如此方式结合在一起，使得其中一个数的每个单位总是对应于另一个数的每个单位时，我们就宣布它们相等。
  \citep[Pt.III Bk.1 \S1]{Hume1740}
\end{quote}

休谟原则被 \citet[\S63]{Frege1884} 带到了当代数理逻辑研究的显著位置，他引用了休谟的这段话，然后（实际上）概述了（我们所说的）休谟原则。

你应该注意到休谟原则与基本定律 V 在结构上的相似性。我们在 \olref[story][blv]{sec} 中将其表述如下：
\[
	\fregeext{x}{F(x)} = \fregeext{x}{G(x)} \text{当且仅当 } \lforall[x][(F(x) \liff G(x))].
\]
在此，一些评论和比较或许有所帮助。

理解像休谟原则或基本定律 V 这样的原理有两种方式：\emph{直谓地}或\emph{非直谓地}（回忆 \olref[story][predicative]{sec}）。在基本定律 V 的非直谓解读下，对于每个 $F$，对象 $\fregeext{x}{F(x)}$ 都落在我们表述基本定律 V 本身时所使用的量化论域中。类似地，在休谟原则的非直谓解读下，对于每个 $F$，对象 $\fregenum{x}{F(x)}$ 也落在我们表述休谟原则时所使用的量化论域中。相比之下，在\emph{直谓}的理解下，对象 $\fregeext{x}{F(x)}$ 和 $\fregenum{x}{F(x)}$ 将是来自某个\emph{不同}论域的实体。

现在，如果我们以非直谓的方式解读基本定律 V，它会通过朴素概括（详情参见 \olref[story][blv]{sec}）导致不一致。与朴素概括非常相似，通过以\emph{直谓}的方式解读它，可以使它变得一致。但这可能无法实现我们希望它做到的所有事情。

然而，休谟原则\emph{可以}一致地以非直谓的方式解读。而且，这样解读时，它相当强大。

举例说明：考虑谓词“$x \neq x$”，显然没有东西满足它。休谟原则现在产生一个对象 $\# x( x\neq x)$。我们可以将其视为数 $0$。现在，在\emph{非直谓}的理解下——但\emph{只有在}非直谓的理解下——这个实体 $0$ 落在我们最初的量化论域中。因此，我们可以合理地将休谟原则应用于谓词“$x = 0$”以获得一个对象 $\#x (x = 0)$。我们可以将其视为数 $1$。此外，休谟原则后承 $0 \neq 1$，因为不可能存在一个从不等同于自身的对象到等同于 $0$ 的对象之间的双射（前者数量为零，而后者有一个）。现在，再次采用非直谓的方式，$1$ 落在我们最初的量化论域中。因此，我们可以合理地将休谟原则应用于谓词“$(x = 0 \lor x = 1)$”以获得一个对象 $\#x(x = 0 \lor
x = 1)$。我们可以将其视为数 $2$，并且我们可以证明 $0\neq 2$ 和 $1 \neq 2$，等等。

简而言之，以非直谓方式理解时，休谟原则后承存在\emph{无穷多个对象}。这促使\emph{新弗雷格主义逻辑主义者}将休谟原则作为算术的基础。

然而，弗雷格\emph{本人}并没有将休谟原则作为他的算术基础。相反，弗雷格基于一个显式定义证明了休谟原则：$\fregenum{x}{F(x)}$ 被定义为概念 $F \sim \Phi$ 的外延。用现代术语来说，我们可能会尝试将其表述为 $\fregenum{x}{F(x)} = \Setabs{G}{F \sim G}$；但这将把我们重新拉回朴素概括的问题之中。

\end{document}